\capituloPregunta{¿Cómo adaptar el diseño experimental a procesos reales?}

\section{Motivación}

En un ejemplo ideal, todos los tratamientos tienen el mismo número de repeticiones, las condiciones se controlan perfectamente y los datos llegan completos. En la práctica, los procesos reales rara vez son tan ordenados.

\section{Caso integrado: proceso manual no balanceado}

El material del curso incluye un caso de optimización de un proceso manual mediante un diseño factorial no balanceado. El problema consiste en estudiar cómo ciertas condiciones del procedimiento afectan el tiempo y la calidad del resultado.

Este caso es importante porque muestra una dificultad frecuente: el diseño experimental debe adaptarse a restricciones reales sin perder rigor.

\section{Qué cambia frente al diseño ideal}

\begin{itemize}
  \item Puede haber repeticiones desiguales entre tratamientos.
  \item Algunas combinaciones pueden ser difíciles de ejecutar.
  \item La habilidad del operador puede introducir variación.
  \item La medición de calidad puede tener un componente subjetivo.
  \item La interpretación debe declarar límites.
\end{itemize}

\begin{cuidado}
Un diseño no balanceado no invalida automáticamente un estudio, pero obliga a ser más cuidadoso con el análisis y con las conclusiones.
\end{cuidado}

\section{Actividad}

Seleccione un proceso manual: dibujar una figura, cortar un material, ensamblar una pieza o preparar una mezcla. Defina tres factores operativos y discuta qué condiciones impedirían obtener un diseño perfectamente balanceado.

\section{Conexión}

Cuando el número de factores crece, probar todas las combinaciones puede volverse costoso. Ahí aparecen los diseños reducidos.
